\documentclass[11pt,a4paper]{article}

% ---------- packages ----------
\usepackage[utf8]{inputenc}
\usepackage[T1]{fontenc}
\usepackage{lmodern}
\usepackage{microtype}
\usepackage[a4paper,margin=1in]{geometry}
\usepackage{amsmath,amssymb,amsthm}
\usepackage{mathtools}
\usepackage{booktabs}
\usepackage{array}
\usepackage{enumitem}
\usepackage{graphicx}
\usepackage{longtable}
\usepackage{xcolor}
\definecolor{accent}{HTML}{1F3A5F}
\usepackage[colorlinks=true,linkcolor=accent,citecolor=accent,urlcolor=accent]{hyperref}
\usepackage{sectsty}
\usepackage{fancyhdr}
\usepackage{setspace}
\onehalfspacing

% ---------- header/footer ----------
\pagestyle{fancy}
\fancyhf{}
\rhead{\small\itshape Limits to Arbitrage --- Literature Review}
\lhead{\small BS Project (EC 495/497)}
\cfoot{\thepage}
\renewcommand{\headrulewidth}{0.4pt}

% ---------- section styling ----------
\partfont{\color{accent}}
\sectionfont{\color{accent}}
\subsectionfont{\color{accent}}

% ---------- math shortcuts ----------
\newcommand{\E}{\mathbb{E}}
\newcommand{\Var}{\operatorname{Var}}
\newcommand{\Cov}{\operatorname{Cov}}
\newcommand{\Corr}{\operatorname{Corr}}
\newcommand{\R}{\mathbb{R}}
\newcommand{\dd}{\mathrm{d}}
\newcommand{\1}{\mathbf{1}}
\DeclarePairedDelimiter{\abs}{\lvert}{\rvert}
\newcommand{\rupee}{\ensuremath{\mathrm{Rs}\,}}
\newtheorem*{takeaway}{Takeaway for our project}

% ---------- title ----------
\title{\vspace{-1.2cm}\color{accent}\bfseries Limits to Arbitrage in the Indian Equity Market\\[0.3cm]
\large\mdseries\color{black} A Mathematically Detailed Literature Review\\
\normalsize\itshape Theory, measurement, and evidence --- assembled to underpin a formal model of the project}
\author{Prince Kumar \quad(Roll No.\ 23B3032)\\ \small BS Economics, IIT Bombay \\[0.2cm]
\small Project guide: Prof.\ Aditi Chaubal \quad\textbar\quad Courses: EC~495 (Stage~I) / EC~497 (Stage~II)}
\date{September 2026}

\begin{document}
\maketitle
\thispagestyle{empty}

\begin{abstract}
\noindent This document is a deep, equation-level review of the twenty-one references assembled for the BS Project on \emph{limits to arbitrage in the Indian equity market}. Its purpose is deliberately narrow and deliberately demanding: to reconstruct, in full mathematical detail, the theoretical models, the mispricing-measurement definitions, and the econometric estimators that these papers contain, so that they can serve as the raw material for the next stage of the project --- a self-contained theoretical/mathematical model of what the project sets out to establish, built \emph{before} any empirical work in order to guard against confirmation bias and data-snooping. Each paper is treated in its own right (question, formal model with derivations, data and method, results) and then mapped explicitly onto our research design. The review is organised into six parts: (I) the theory of limited arbitrage; (II) how mispricing is measured; (III) the econometrics of the no-arbitrage band; (IV) evidence from the Indian market and the institutional mechanism we exploit; (V) the behavioural/retail demand side; and (VI) a synthesis that states precisely what a theoretical model of our project must reproduce. Throughout, the recurring object is the \emph{no-arbitrage band} --- the interval of mispricing inside which arbitrage does not pay --- and the recurring theme is that India's mechanical F\&O ``ban'' on single-stock derivatives is a rare, dated, exogenous shock to the width of that band.
\end{abstract}

\vspace{0.4cm}
\begin{center}
\fbox{\parbox{0.92\textwidth}{\small
\textbf{Our problem statement (restated for reference).}
Standard theory says a mispriced security is traded back to fair value almost at once. In practice, price errors in Indian stocks survive for days or weeks. The project asks: \emph{how strongly do arbitrage constraints bind in Indian equities, and how much of the observed, persistent mispricing can be attributed to them?} The identifying variation is India's Market-Wide Position Limit (MWPL) ``F\&O ban'': once market-wide open interest in a single stock crosses 95\% of its MWPL, no fresh derivative positions may be opened. Because single-stock futures are the primary arbitrage channel in India (intraday cash short-selling is allowed, but naked overnight shorting is barred and the securities-lending market is thin), the ban is a sharp, mechanical, repeated contraction of arbitrage capacity.}}
\end{center}

\newpage
\tableofcontents
\newpage

% ---------- how to read ----------
\section*{How to read this document}
\addcontentsline{toc}{section}{How to read this document}
Every paper is written up so that a reader who has never opened it can follow both \emph{what it does} and \emph{the mathematics of how}. Sections follow a common template: the question and contribution; the formal model or measurement definition, with derivations; data and method; results; and a boxed \textit{Takeaway for our project}. Notation is kept consistent across papers wherever possible --- in particular $F$ and $S$ (or $P$) denote futures and spot prices, $z$ or $b$ the basis/mispricing, $c$ a per-unit friction, and $u$ a non-fundamental demand shock. Cross-references between papers are frequent and deliberate: the value of the review is in how the strands fit together, which Part~VI makes explicit. A consolidated reference list appears at the end.

\part{The Theory of Limited Arbitrage}
\label{part:theory}

The three papers in this part answer a single question in three complementary ways: \emph{why does a mispricing survive at all?} Classical theory (Fama, Ross) assumes a continuum of small, diversified, infinitely-lived arbitrageurs who compete mispricing away to zero. Each model below breaks one of those assumptions --- limited/withdrawable capital (Shleifer--Vishny), short horizons facing self-created risk (De~Long et al.), and explicit position/financing frictions (Gromb--Vayanos) --- and each delivers the same qualitative object: an interval of mispricing, a \emph{no-arbitrage band}, inside which rational arbitrage does not operate. Together they are the theoretical justification for everything measured in Parts~II--IV.

% =====================================================================
\section{Shleifer \& Vishny (1997): Performance-Based Arbitrage}
\label{sec:sv}

\subsection*{Question and contribution}
Standard models predict that the \emph{larger} a mispricing, the \emph{more} arbitrage capital rushes in to correct it, so arbitrage is most powerful exactly when it is most needed. Shleifer and Vishny (SV) overturn this by taking seriously that real arbitrage is done by \emph{specialists trading other people's money}. The separation of ``brains'' (the arbitrageur's knowledge) from ``resources'' (outside investors' capital) creates an agency problem: capital is allocated on \emph{past} performance, which is worst precisely when \emph{future} expected returns are highest. The result is that arbitrage can be weakest in extreme mispricings --- and can even amplify them.

\subsection*{The model}
Consider one asset with fundamental value $V$, revealed at $t=3$. There are three periods $t=1,2,3$ and three types of agent: noise traders, arbitrageurs, and the outside investors who fund them. In periods $t=1,2$ pessimistic noise traders exert an aggregate ``pessimism shock'' $S_t$ that depresses demand; their demand for the asset is
\begin{equation}
Q^{N}(t) \;=\; \frac{V - S_t}{p_t},
\end{equation}
where $p_t$ is the price. The first shock $S_1$ is known at $t=1$; the second, $S_2$, is uncertain:
\begin{equation}
S_2 = \begin{cases} S \;(>S_1) & \text{with probability } q \quad(\text{misperception deepens}),\\[2pt] 0 & \text{with probability } 1-q \quad(\text{noise corrects: } p_2=V).\end{cases}
\end{equation}
Risk-neutral arbitrageurs hold cumulative funds $F_t$ and invest $D_t \le F_t$ in the asset at $t=1$ (keeping $F_1-D_1$ in cash as an option to buy more if the asset gets cheaper). With unit asset supply, market clearing $Q^N(t)+Q^A(t)=1$ (with $Q^A=D_t/p_t$ or $F_t/p_t$) gives the prices
\begin{align}
p_1 &= V - S_1 + D_1, \label{eq:sv-p1}\\
p_2 &= V - S_2 + F_2. \label{eq:sv-p2}
\end{align}
The maintained assumption is that arbitrage capital is too small to remove the mispricing, $F_t < S_t$, so prices stay below $V$ whenever noise traders are active.

\paragraph{Performance-based arbitrage (PBA).} The engine of the model is how $F_2$ is determined. Outside investors cannot evaluate strategies; they update on realised returns only. If the arbitrageur's gross return between $t=1$ and $t=2$ is $p_2/p_1$ on the invested part, funds next period are
\begin{equation}
F_2 \;=\; G\!\left(D_1\frac{p_2}{p_1} + (F_1 - D_1)\right), \qquad G'>0,\quad G''\le 0, \label{eq:sv-G}
\end{equation}
with $G$ normalised so that performance equal to the benchmark leaves funds unchanged, and with fund-flow sensitivity $G'$ large enough that under-performance triggers withdrawals. This is \emph{PBA}: capital chases past returns. Because all arbitrageurs in a segment hold the same position, they gain or lose funds together --- there is no deep-pocketed outsider to step in.

\subsection*{Solving for equilibrium}
Since managers are paid in proportion to funds managed, they maximise expected terminal funds. If noise corrects at $t=2$ the position is liquidated at $V$ and held in cash; if it deepens, funds $F_2$ then grow by $V/p_2$ to $t=3$. The objective is
\begin{equation}
\E W \;=\; (1-q)\,G\!\left(D_1\frac{V}{p_1} + (F_1-D_1)\right) \;+\; q\,\frac{V}{p_2}\,G\!\left(D_1\frac{p_2}{p_1} + (F_1-D_1)\right). \label{eq:sv-EW}
\end{equation}
The first-order condition in $D_1$ trades off the benefit of being fully invested if the price recovers against the option value of holding cash in case it falls further. Arbitrageurs choose to be \emph{fully invested} ($D_1=F_1$) when the recovery is likely and the initial displacement large --- precisely,
\begin{equation}
(1-q)\,G'\!\left(\tfrac{V}{p_1}-1\right) \;+\; q\,G'\!\left(\tfrac{p_2}{p_1}-1\right)\frac{V}{p_2} \;>\;0 \quad\text{at } D_1=F_1. \label{eq:sv-fullinv}
\end{equation}
Under full investment, $F_2 = G(F_1\,p_2/p_1)$, and substituting into \eqref{eq:sv-p2} gives the second-period price \emph{implicitly}:
\begin{equation}
p_2 \;=\; V - S + G\!\left(F_1\frac{p_2}{p_1}\right). \label{eq:sv-p2imp}
\end{equation}

\subsection*{The key result: negative resiliency}
Differentiate \eqref{eq:sv-p2imp} with respect to the deepening shock $S$ (treating $p_1$ as given):
\begin{equation}
\frac{\dd p_2}{\dd S} \;=\; \frac{-1}{\,1 - G'\dfrac{F_1}{p_1}\,} \;<\; -1, \label{eq:sv-resiliency}
\end{equation}
where the stable equilibrium requires $0 < G'F_1/p_1 < 1$. Equation~\eqref{eq:sv-resiliency} is the paper's punchline. In an efficient market the sensitivity of price to a sentiment shock would be zero; with limited capital it is $-1$ at $t=1$; but under PBA, at $t=2$ the price falls \emph{more than one-for-one} with the shock. When the mispricing worsens, arbitrageurs suffer losses, investors withdraw funds ($F_2$ falls), and the arbitrageurs are forced to \emph{liquidate into the falling market} --- a fire sale that pushes the price still further from $V$. Arbitrage is not merely limited; at the extreme it is destabilising.

\begin{takeaway}
SV is the conceptual charter for the whole project. It says (i) arbitrage capacity is finite and \emph{state-dependent} --- it shrinks in bad states; (ii) mispricing has a \emph{time dimension} (it can widen and persist), so we must model the \emph{speed} of correction, not just its level (Parts~II--III); and (iii) a sudden, external contraction of arbitrage capacity should \emph{widen} mispricing. India's MWPL ban is exactly such a contraction, imposed mechanically and dated to the day. The amplification channel \eqref{eq:sv-resiliency} also warns us that ban effects may be non-linear and largest in already-stressed names.
\end{takeaway}

% =====================================================================
\section{De Long, Shleifer, Summers \& Waldmann (1990): Noise-Trader Risk}
\label{sec:dssw}

\subsection*{Question and contribution}
Friedman and Fama argued that irrational ``noise traders'' cannot move prices for long: rational arbitrageurs trade against them, profit, and drive them extinct. DSSW show this argument fails once arbitrageurs are \emph{risk-averse with short horizons}. The very unpredictability of noise traders' future beliefs is a risk --- \emph{noise-trader risk} --- that deters arbitrage. Prices then diverge from fundamentals \emph{even with no fundamental risk at all}, and noise traders can earn \emph{higher} average returns by bearing the risk they themselves create. ``Noise traders create their own space.''

\subsection*{The model}
An overlapping-generations economy: two-period-lived agents, no first-period consumption, so investable resources are exogenous. Two assets pay the identical dividend $r$: a safe asset $s$ in perfectly elastic supply (price fixed at $1$, hence $r$ is the riskless rate), and a risky asset $u$ in fixed unit supply with price $p_t$. If priced at fundamentals both would equal $1$ in every period; noise traders break this. A fraction $\mu$ of agents are noise traders, $1-\mu$ are sophisticated. Every agent has CARA utility over end-of-life wealth $w$,
\begin{equation}
U(w) = -\exp(-2\gamma w),
\end{equation}
with absolute risk aversion $2\gamma$. Under normality, maximising $\E U$ is equivalent to maximising $\bar w - \gamma\,\sigma_w^2$. Noise traders misperceive the expected \emph{next-period price} of $u$ by an i.i.d.\ normal amount
\begin{equation}
\rho_t \sim \mathcal N\!\left(\rho^\ast,\ \sigma_\rho^2\right),
\end{equation}
where $\rho^\ast$ measures average ``bullishness'' and $\sigma_\rho^2$ the variability of sentiment.

\paragraph{Demands.} Let ${}_t\sigma^2_{p_{t+1}}\equiv \E_t[(p_{t+1}-\E_t p_{t+1})^2]$ be the one-period-ahead price variance. Mean--variance optimisation gives holdings of the risky asset:
\begin{align}
\lambda^{i}_t &= \frac{r + {}_t\E p_{t+1} - (1+r)p_t}{2\gamma\, {}_t\sigma^2_{p_{t+1}}} && \text{(sophisticated)}, \label{eq:dssw-lami}\\[4pt]
\lambda^{n}_t &= \frac{r + {}_t\E p_{t+1} - (1+r)p_t}{2\gamma\, {}_t\sigma^2_{p_{t+1}}} \;+\; \frac{\rho_t}{2\gamma\, {}_t\sigma^2_{p_{t+1}}} && \text{(noise traders)}. \label{eq:dssw-lamn}
\end{align}
Demand is (perceived excess return)/(risk aversion $\times$ variance); the only difference is the extra term $\rho_t/(2\gamma\,{}_t\sigma^2_{p_{t+1}})$ by which optimistic noise traders over-hold. Crucially, the \emph{variance in the denominator is generated endogenously by noise trading itself}: even sophisticated investors cannot take unbounded positions, because the price at which they must sell when old depends on next period's (uncertain) noise-trader mood. This is what prevents arbitrage from enforcing the law of one price.

\subsection*{The pricing function}
Old agents sell their unit; young agents' demands must sum to the unit supply. Imposing this and looking for a steady state (the unconditional distribution of $p_{t+1}$ equals that of $p_t$) yields the central closed form,
\begin{equation}
\boxed{\;p_t \;=\; 1 \;+\; \underbrace{\frac{\mu(\rho_t-\rho^\ast)}{1+r}}_{\text{(a) sentiment fluctuation}} \;+\; \underbrace{\frac{\mu\rho^\ast}{r}}_{\text{(b) mean mispricing}} \;-\; \underbrace{\frac{(2\gamma)}{r}\,\frac{\mu^2\sigma_\rho^2}{(1+r)^2}}_{\text{(c) noise-trader-risk premium}}\;}
\label{eq:dssw-price}
\end{equation}
with the (constant) one-step-ahead price variance
\begin{equation}
{}_t\sigma^2_{p_{t+1}} \;=\; \frac{\mu^2\sigma_\rho^2}{(1+r)^2}. \label{eq:dssw-var}
\end{equation}
Read \eqref{eq:dssw-price} term by term. (a) The price moves one-for-one (scaled by $\mu/(1+r)$) with \emph{current} sentiment $\rho_t$ --- excess volatility unrelated to fundamentals. (b) If noise traders are bullish on average ($\rho^\ast>0$) the price is permanently \emph{above} fundamental value. (c) Regardless of the sign of $\rho^\ast$, the term $-(2\gamma/r)\,\mu^2\sigma_\rho^2/(1+r)^2$ is a strict \emph{discount}: the mere presence of sentiment \emph{risk} lowers the price. It is this term that survives even when average sentiment is neutral, and that is the formal statement of ``limits to arbitrage without fundamental risk.''

\subsection*{Relative returns: noise traders can win}
Because bullish noise traders hold more of the risky asset (which carries the risk premium (c)), they can earn a \emph{higher} expected return than the sophisticates who lean against them. DSSW show the return differential is non-monotone (inverted-U) in $\rho^\ast$: a little bullishness earns the risk premium and beats the arbitrageurs, but too much bullishness makes noise traders buy at prices so inflated that they lose. This is the mechanism behind the SEBI finding (Part~V) that the \emph{aggregate} of retail derivative traders loses to proprietary and foreign desks: crowding is self-defeating past a point, even as it moves prices.

\begin{takeaway}
DSSW supplies the second leg of the band (SV supplied the capital leg): even absent any capital-withdrawal channel, \emph{risk} created by unpredictable order flow keeps a finite wedge open. For us, three things carry over. First, the discount term (c) shows mispricing scales with the \emph{variance} of non-fundamental demand $\sigma_\rho^2$ --- a rationale for putting idiosyncratic volatility \emph{into} our constraint index. Second, the model is explicitly about \emph{marketwide} (non-diversifiable) sentiment, which is why our constraint index will have a systematic (PCA) component. Third, it is the theoretical bridge to the retail/policy chapter: retail crowding is noise-trader demand, and the ban is triggered when that demand piles into a single name.
\end{takeaway}

% =====================================================================
\section{Gromb \& Vayanos (2010): A Unified Model of Arbitrage Costs}
\label{sec:gv}

\subsection*{Question and contribution}
Gromb and Vayanos (GV) survey the field and, importantly for us, \emph{nest} its mechanisms in one tractable model, cataloguing the costs an arbitrageur faces: (i) risk (fundamental and non-fundamental), (ii) short-selling costs, (iii) leverage/margin constraints, and (iv) constraints on equity capital. The model shows how each cost maps into a \emph{price wedge}, and --- the single most useful result for us --- how a short-selling cost makes that wedge \emph{one-sided}, i.e.\ produces an \emph{asymmetric} no-arbitrage band.

\subsection*{Baseline: the price impact of a demand shock}
Two risky assets $A,B$ trade in Period~1 and pay off in Period~2; the riskless rate is $0$. Payoffs $d_A,d_B$ are jointly normal with means $\bar d_i$, standard deviations $\sigma_i$, correlation $\rho$. Asset $B$ is priced at its expected payoff, $p_B=\bar d_B$; $p_A$ is endogenous. Outside investors demand $u$ shares of $A$ inelastically --- $u$ is the \emph{demand shock}. Competitive arbitrageurs with CARA utility choose positions $x_A,x_B$ to maximise
\begin{equation}
-\,\E_1\!\left[\exp(-\alpha W_2)\right],\qquad W_2 = W_1 + x_A(d_A-p_A) + x_B(d_B-p_B). \label{eq:gv-obj}
\end{equation}
Using normality and $p_B=\bar d_B$, this reduces to the mean--variance problem
\begin{equation}
\max_{x_A,x_B}\; x_A(\bar d_A - p_A) - \frac{\alpha}{2}\Big(x_A^2\sigma_A^2 + x_B^2\sigma_B^2 + 2x_Ax_B\rho\sigma_A\sigma_B\Big). \label{eq:gv-mv}
\end{equation}
The optimal hedge is $x_B = -(\rho\sigma_A/\sigma_B)\,x_A$, and the optimal position in $A$ is
\begin{equation}
x_A = \frac{\bar d_A - p_A}{\alpha\sigma_A^2(1-\rho^2)}. \label{eq:gv-xA}
\end{equation}
Market clearing $x_A + u = 0$ (asset $A$ in zero net supply) gives the equilibrium price
\begin{equation}
\boxed{\;p_A \;=\; \bar d_A + \alpha\,\sigma_A^2\,(1-\rho^2)\,u\;} \label{eq:gv-price}
\end{equation}
So a demand shock moves the price, and its impact is \emph{larger} when arbitrageurs are more risk-averse ($\alpha$), when the asset is riskier ($\sigma_A$), and when it is harder to hedge (small $|\rho|$; more idiosyncratic risk, fewer substitutes). If $A$ and $B$ are perfect substitutes ($\rho=1$) the wedge vanishes: arbitrage is riskless and perfect. This equation is the ``how big is the band'' benchmark; everything else adds terms to it. (Under the ``intertemporal arbitrage'' reading, $A$ and $B$ are the same asset at two dates and $\sigma_A^2(1-\rho^2)$ becomes the risk borne while waiting for capital to arrive.)

\subsection*{Non-fundamental risk feeds back}
When the shock $u$ is itself random (introduced in a Period~0), \eqref{eq:gv-price} implies asset $A$'s price volatility rises to
\begin{equation}
\sigma_A\sqrt{\,1 + \alpha^2\sigma_A^2(1-\rho^2)^2\sigma_u^2\,}, \label{eq:gv-vol}
\end{equation}
and its correlation with $B$ falls. Demand-shock risk thus \emph{amplifies itself} (as in DSSW): the possibility of future shocks makes short-horizon arbitrageurs less willing to absorb today's shock.

\subsection*{Short-selling cost: the origin of asymmetry}
Now let shorting $A$ cost $c$ per share. The objective \eqref{eq:gv-mv} gains a term $-\,c\,\abs{x_A}\,\1\{x_A<0\}$, and the equilibrium price becomes \emph{one-sided}:
\begin{equation}
p_A = \begin{cases} \bar d_A + \alpha\sigma_A^2(1-\rho^2)\,u, & u \le 0 \quad(\text{arbitrageurs long}),\\[4pt] \bar d_A + \alpha\sigma_A^2(1-\rho^2)\,u + c, & u > 0 \quad(\text{arbitrageurs short}). \end{cases} \label{eq:gv-asym}
\end{equation}
The short-selling cost bites \emph{only} when arbitrageurs must short (positive demand shock), lifting the price by exactly $c$. Setting $\rho=1$ (identical payoffs) isolates a pure law-of-one-price violation:
\begin{equation}
p_A = \bar d + c \quad\text{when } u>0, \qquad p_A=\bar d \quad\text{when } u\le 0.
\end{equation}
An identical-payoff asset is over-priced by the cost of shorting it, and not at all under-priced.

\subsection*{Leverage/margin and equity-capital constraints}
When arbitrageurs face a margin requirement $\sum_i m_i\abs{x_i}\le W$, the constraint binds in states where the position is large; its Lagrange multiplier acts as an \emph{additional, endogenous holding cost} that widens the wedge exactly when the shock (and hence the opportunity) is largest. Equity-capital constraints reintroduce the SV channel: a loss shrinks $W$, tightens the constraint, forces de-leveraging, and can trigger the fire-sale amplification of \S\ref{sec:sv}. GV's synthesis is that all four costs are variations on one theme --- each adds a term to \eqref{eq:gv-price} that is increasing in the size of the required position and in how binding the friction is.

\begin{takeaway}
GV is the model our project most directly instantiates, for two reasons. (1) It gives the exact functional shape of the object we estimate: mispricing $=$ (risk-bearing term) $+$ (friction wedge), with the wedge \emph{one-sided} under short-sale-type constraints --- which predicts and explains the \emph{asymmetric} bands found for India (Aggarwal; Ofek; \S\ref{sec:aggarwal},~\ref{sec:ofek}). (2) India's MWPL ban is a \emph{hard position/leverage constraint} of exactly GV type (iii): once binding, no fresh positions --- the shadow price of the constraint jumps, and the wedge widens. Our constraint index is, in GV's language, an empirical proxy for the combined shadow cost of costs (ii)--(iv); our causal tests ask whether an exogenous increase in that shadow cost (a ban) moves the wedge as \eqref{eq:gv-asym} predicts.
\end{takeaway}

\subsection{Part~I synthesis: the anatomy of a no-arbitrage band}
The three models compose into one picture. Write the observed price of a security as its fundamental value plus a wedge:
\begin{equation}
p = \underbrace{V}_{\text{fundamental}} \;+\; \underbrace{\theta(u,\sigma^2,\alpha)}_{\substack{\text{risk-bearing term}\\ \text{(DSSW, GV baseline)}}} \;\pm\; \underbrace{c(\text{constraints})}_{\substack{\text{one-sided friction wedge}\\ \text{(GV short-sale/margin)}}} \;+\; \underbrace{\text{amplification}}_{\substack{\text{state-dependent, PBA}\\ \text{(SV)}}}.
\end{equation}
This decomposition motivates the empirical programme exactly: Part~II measures the wedge (the basis, put--call-parity deviations, PEAD); Part~III estimates the \emph{band} --- the region where the wedge persists because it is smaller than the round-trip friction --- and its asymmetry; and Part~IV asks whether an exogenous shock to $c$ (the MWPL ban and the 2024--26 reforms) shifts the band as theory predicts. The persistent \emph{negative} basis documented for India (Part~IV) is the empirical signature of the one-sided wedge in \eqref{eq:gv-asym}, since the binding constraint in India is on the short/negative-basis side.

\part{Measuring Mispricing}
\label{part:measurement}

Part~I established \emph{that} a wedge exists. To test anything we must \emph{measure} it. This part sets out the three mispricing measures the project will construct, each of which is a way of pricing ``the same payoff twice'' and reading the gap: the cost-of-carry futures basis, put--call-parity deviations, and post-earnings-announcement drift. It closes with Stambaugh--Yu--Yuan's composite anomaly score, which is both a mispricing measure and the template for our cross-sectional test.

% =====================================================================
\section{The cost-of-carry basis}
\label{sec:basis}
Every futures--spot paper in the bibliography (Dwyer et al.; Nandan et al.; Tse; Aggarwal; Roy \& Chakraborty) rests on the same no-arbitrage identity. For a stock with spot price $S_t$, continuously-compounded interest rate $r$, dividend yield $\delta$, and time to expiry $\tau=T-t$, the frictionless fair futures price is
\begin{equation}
F_t^{\ast} = S_t\, e^{(r-\delta)\tau}. \label{eq:coc}
\end{equation}
The \emph{mispricing} (or basis error) is the deviation of the observed futures price $F_t$ from \eqref{eq:coc}. Two equivalent normalisations appear in the literature:
\begin{align}
z_t &= \frac{F_t - S_t\,e^{(r-\delta)\tau}}{S_t} && \text{(proportional; Aggarwal)}, \label{eq:basis-prop}\\
z_t &= f_t - \big(s_t + (r-\delta)\tau\big),\qquad f_t=\ln F_t,\ s_t=\ln S_t && \text{(log; Dwyer, Tse)}. \label{eq:basis-log}
\end{align}
The log form makes the econometrics clean: if $f_t$ and $s_t$ are each $I(1)$, then \eqref{eq:coc} implies they are \emph{cointegrated with unit coefficient}, and $z_t$ --- the error-correction term --- is stationary. That single fact is what licenses the threshold error-correction models of Part~III. Positive $z_t$ (futures rich / stock cheap) is corrected by \emph{cash-and-carry} arbitrage (buy spot, sell futures); negative $z_t$ (futures cheap / stock rich) requires \emph{reverse cash-and-carry} (short spot, buy futures) --- the leg that India's constraints make expensive, which is why the observed basis is asymmetric (Part~IV).

\begin{takeaway}
Equation~\eqref{eq:basis-prop} is our primary mispricing measure (Chapter~2 of the thesis). It is computable for every F\&O stock every day from public settlement data, and its sign lines up mechanically with the constrained leg. This is the empirical counterpart of the GV wedge in \eqref{eq:gv-asym}.
\end{takeaway}

% =====================================================================
\section{Ofek, Richardson \& Whitelaw (2004): Put--Call-Parity Deviations}
\label{sec:ofek}

\subsection*{Question and contribution}
A second way to price a stock twice is through its options. Put--call parity (PCP) ties the stock price to a synthetic stock built from a call, a put and a bond. If shorting the stock is costly, the arbitrage that enforces PCP breaks on one side only. Ofek--Richardson--Whitelaw (ORW) measure PCP violations across the whole US single-name option universe, tie them directly to a market-based measure of shorting cost (the rebate-rate spread), and show the violations \emph{predict} future stock returns --- the cleanest available demonstration that short-sale constraints sustain over-pricing.

\subsection*{The measurement}
For European options on a non-dividend stock with strike $K$, maturity $T$, call $C$ and put $P$, PCP is the identity
\begin{equation}
S = \mathrm{PV}(K) + C - P, \qquad \mathrm{PV}(K)=K e^{-rT}. \label{eq:pcp-eu}
\end{equation}
US single-name options are American; early exercise makes the put worth more, so parity becomes an inequality with an early-exercise premium $\mathrm{EEP}\ge 0$:
\begin{equation}
S \;\ge\; \mathrm{PV}(K) + C - P + \mathrm{EEP}. \label{eq:pcp-am}
\end{equation}
Define the \emph{option-implied stock value} and the proportional deviation
\begin{equation}
S^{\ast} \equiv \mathrm{PV}(K) + C - P + \mathrm{EEP}, \qquad \text{deviation} = \frac{S - S^{\ast}}{S}. \label{eq:pcp-dev}
\end{equation}
A positive deviation ($S>S^\ast$) means the stock is expensive relative to its options --- exactly the situation a short-seller would exploit if shorting were free. ORW proxy the difficulty of shorting by the \emph{rebate-rate spread}: the gap between the normal (``cold'') rebate rate on short-sale cash collateral and the (lower) rate on a hard-to-borrow stock. A larger spread $=$ costlier/more-binding short constraint.

\subsection*{Results}
Three findings, all of which we expect --- amplified --- in India:
\begin{enumerate}[leftmargin=1.4em,itemsep=2pt]
\item \textbf{Asymmetry.} After accounting for shorting costs, $13.6\%$ of stocks lie \emph{above} the upper no-arbitrage bound (over-priced) versus only $4.4\%$ below the lower bound; the mean deviation among violators is $2.71\%$. Violations point in the direction the short-sale constraint predicts --- the empirical image of GV's one-sided wedge \eqref{eq:gv-asym}.
\item \textbf{Dose--response.} Both the probability and the size of violations rise with the rebate-rate spread; a one-standard-deviation fall in the rebate spread widens the deviation by $\approx 0.67\%$, robust to option/equity liquidity and transaction costs.
\item \textbf{Predictability.} Larger PCP violations (high implied over-pricing) forecast \emph{lower} future stock returns; long-industry/short-flagged-stock portfolios earn cumulative abnormal returns as high as $\sim 65\%$ over the sample --- mispricing that persists because it cannot be arbitraged.
\end{enumerate}

\begin{takeaway}
ORW gives us the second mispricing measure \eqref{eq:pcp-dev} (Chapter~2) and a discipline: restrict PCP tests to near-the-money strikes, where Indian single-stock option quotes are reliable. Their asymmetry and predictability results are hypotheses H1--H2 in a US mirror; India's near-absolute overnight short constraint means we expect the over-pricing side to be \emph{more} constrained than ORW found. The rebate-rate spread has no clean Indian analogue (SLB is thin and lightly used), which is precisely why we substitute the futures basis and the MWPL ban as our constraint signals.
\end{takeaway}

% =====================================================================
\section{Post-Earnings-Announcement Drift (Harshita, Singh \& Yadav, 2018)}
\label{sec:pead}

\subsection*{Question and contribution}
The third mispricing measure is dynamic rather than cross-sectional. Post-earnings-announcement drift (PEAD) is the tendency of a stock's price to keep drifting in the direction of an earnings surprise for weeks after the announcement --- a violation of semi-strong efficiency, because the information is already public. Harshita et al.\ document robust PEAD in India (2002--2017), which validates it as a usable mispricing signal for our setting.

\subsection*{The measurement}
The surprise is captured by \emph{standardised unexpected earnings} (SUE). With a seasonal-random-walk earnings expectation,
\begin{equation}
\mathrm{SUE}_{i,q} = \frac{X_{i,q} - X_{i,q-4}}{\sigma_{i,q}}, \label{eq:sue}
\end{equation}
where $X_{i,q}$ is firm $i$'s earnings in quarter $q$, $X_{i,q-4}$ the same quarter a year earlier, and $\sigma_{i,q}$ the standard deviation of the surprise over a trailing window. Stocks are sorted into SUE deciles; the drift is the post-event cumulative abnormal return
\begin{equation}
\mathrm{CAR}_i(t_1,t_2) = \sum_{t=t_1}^{t_2}\big(R_{i,t} - R_{i,t}^{\text{benchmark}}\big), \label{eq:car}
\end{equation}
measured over, e.g., $[+1,+60]$ trading days after the announcement. PEAD is present if CAR is monotone and increasing in the SUE decile. Harshita et al.\ confirm statistically significant drift, robust across sub-periods and after controlling for beta, market capitalisation, price-to-book, illiquidity and idiosyncratic volatility, using cross-sectional Fama--MacBeth regressions (see \S\ref{sec:syy} for the estimator).

\begin{takeaway}
PEAD is our third mispricing measure (Chapter~2) and the one with the sharpest identification: earnings dates are exogenous and precisely timed. Hypothesis H2 is that drift is \emph{larger} for constrained stocks and \emph{largest on days a stock is in the ban}, because the trades that would close the gap (often on the short side) are exactly the ones the ban blocks. That the Indian PEAD paper already controls for idiosyncratic volatility tells us our constraint index must add \emph{incremental} explanatory power beyond IVOL to be a contribution.
\end{takeaway}

% =====================================================================
\section{Stambaugh, Yu \& Yuan (2015): A Composite Mispricing Score}
\label{sec:syy}

\subsection*{Question and contribution}
Stambaugh--Yu--Yuan (SYY) resolve the ``idiosyncratic volatility puzzle'' (why high-IVOL stocks earn low average returns) with a mechanism we reuse wholesale: \emph{arbitrage asymmetry}. Shorting is harder than buying, so over-pricing is corrected less than under-pricing. Combined with a proxy for mispricing, this predicts that IVOL --- which measures how risky arbitrage is --- should relate \emph{negatively} to returns among over-priced stocks and \emph{positively} among under-priced ones, with the negative side stronger. Their construction of the mispricing proxy and their IVOL/arbitrage-risk logic are the blueprint for our Chapter~5.

\subsection*{The composite mispricing measure}
No single anomaly cleanly measures mispricing, so SYY \emph{average across anomalies}. For each of $11$ return anomalies that survive Fama--French three-factor adjustment, they rank every stock into a percentile
\begin{equation}
\mathrm{rank}_{i,j,t}\in\{1,\dots,100\}, \qquad \text{higher} = \text{more over-priced by anomaly } j,
\end{equation}
and define the stock's mispricing score as the average rank across the anomalies it is scored on:
\begin{equation}
M_{i,t} = \frac{1}{J_{i,t}}\sum_{j=1}^{J_{i,t}} \mathrm{rank}_{i,j,t}. \label{eq:syy-M}
\end{equation}
A high $M$ flags relative over-pricing. The eleven anomalies span financial distress (failure probability and Ohlson's O-score), net stock issuance, composite equity issuance, total accruals, net operating assets, momentum, gross profitability, asset growth, return on assets, and investment-to-assets. Idiosyncratic volatility is measured as the standard deviation of residuals from a Fama--French three-factor regression over a trailing month:
\begin{equation}
r_{i,t} - r_{f,t} = \alpha_i + \beta_i^{\text{MKT}}\mathrm{MKT}_t + \beta_i^{\text{SMB}}\mathrm{SMB}_t + \beta_i^{\text{HML}}\mathrm{HML}_t + \varepsilon_{i,t},\qquad \mathrm{IVOL}_{i}=\mathrm{sd}(\hat\varepsilon_i). \label{eq:ivol}
\end{equation}

\subsection*{The result and the estimator}
Double-sorting on $M$ and IVOL, SYY find the IVOL--return relation is significantly negative among the most over-priced stocks, positive among the most under-priced, and stronger on the over-priced side --- and that these effects strengthen with the Baker--Wurgler investor-sentiment index in the predicted directions. Their cross-sectional workhorse is the \emph{Fama--MacBeth} two-pass regression: in each month $t$, run the cross-sectional regression
\begin{equation}
r_{i,t+1} = \gamma_{0,t} + \gamma_{1,t}\,\mathrm{IVOL}_{i,t} + \gamma_{2,t}'\,\mathbf{X}_{i,t} + u_{i,t+1}, \label{eq:fm}
\end{equation}
then average the slopes over time, $\hat\gamma_k=\frac1T\sum_t\hat\gamma_{k,t}$, with $t$-statistics from the time series of $\{\hat\gamma_{k,t}\}$ (Newey--West corrected). The sign and significance of $\hat\gamma_1$ within mispricing groups is the test.

\begin{takeaway}
SYY hands us three tools. (1) The composite score \eqref{eq:syy-M} is exactly how we will turn our $\sim$45-signal library into a single mispricing sort. (2) The Fama--MacBeth regression \eqref{eq:fm} is our Chapter~5 estimator; our twist is to interact IVOL / signal returns with the \emph{constraint index}, testing H5 (signals pay more where constraints bind). (3) The arbitrage-asymmetry logic makes a sharp Indian prediction: with the short side near-shut, the over-pricing/negative-basis persistence should be even stronger than in the US --- tying SYY back to the one-sided GV wedge \eqref{eq:gv-asym} and to the negative Indian basis of Part~IV.
\end{takeaway}

\part{The Econometrics of the No-Arbitrage Band}
\label{part:econometrics}

Part~I says a band exists; Part~II measures the deviation $z_t$ that lives inside it. This part is the machinery that \emph{estimates the band itself} --- the width of the interval in which $z_t$ does not mean-revert because it is smaller than the round-trip friction. The unifying idea is \emph{non-linear, regime-dependent} adjustment: $z_t$ behaves like a (near) random walk inside the band and mean-reverts sharply outside it. We build up from the cointegration foundation to the estimator the project will actually use, Ters--Urban's three-regime threshold VECM.

% =====================================================================
\section{Foundation: threshold cointegration}
\label{sec:threshold-coint}
From \eqref{eq:basis-log}, log-futures $f_t$ and log-spot $s_t$ are cointegrated with unit vector, so the basis $z_t = f_t - s_t - (r-\delta)\tau$ is stationary and is the \emph{error-correction term}. Balke \& Fomby (1997) formalise the case where adjustment back to equilibrium happens \emph{only} once the deviation is large enough --- \emph{threshold cointegration}. In their band-TAR the equilibrium error follows
\begin{equation}
\Delta z_t = \begin{cases} \alpha_{\text{out}}\,(z_{t-1}-\gamma_{\text{u}}) + \varepsilon_t, & z_{t-1} \ge \gamma_{\text{u}},\\[3pt] \phantom{\alpha_{\text{out}}(}\varepsilon_t, & \gamma_{\ell} < z_{t-1} < \gamma_{\text{u}} \quad(\text{no-arbitrage band}),\\[3pt] \alpha_{\text{out}}\,(z_{t-1}-\gamma_{\ell}) + \varepsilon_t, & z_{t-1} \le \gamma_{\ell}, \end{cases} \label{eq:bandtar}
\end{equation}
with $\alpha_{\text{out}}<0$ (mean reversion outside) and \emph{no} error-correction term inside $(\gamma_\ell,\gamma_{\text{u}})$. The band $[\gamma_\ell,\gamma_{\text{u}}]$ is the object of interest, and its half-width is an estimate of the round-trip arbitrage cost. Everything below is a version of \eqref{eq:bandtar}: univariate (TAR), smooth (STAR), or bivariate (TVECM).

% =====================================================================
\section{The threshold autoregression: Yadav--Pope--Paudyal and Dwyer et al.\ (1996)}
\label{sec:tar}
Yadav, Pope \& Paudyal (1994) proposed, and Dwyer, Locke \& Yu (1996) developed, the \emph{self-exciting three-regime TAR} for the futures--cash basis. Regimes are switched by the lagged basis $z_{t-d}$ relative to two thresholds $C_1<C_2$:
\begin{equation}
z_t = \phi_0^{(j)} + \sum_{k=1}^{p}\phi_k^{(j)}\,z_{t-k} + \varepsilon_t^{(j)}, \qquad
j = \begin{cases} 1 & z_{t-d}\le C_1 \ (\text{lower}),\\ 2 & C_1<z_{t-d}<C_2 \ (\text{middle}),\\ 3 & z_{t-d}\ge C_2 \ (\text{upper}), \end{cases} \label{eq:tar}
\end{equation}
with $\varepsilon^{(j)}\sim(0,\sigma_j^2)$ i.i.d.\ The \emph{delay} $d$ sets how quickly arbitrageurs react; $p$ is the lag order. Inside the middle regime $z_t$ is (near) non-stationary --- $\sum_k \phi_k^{(2)}\approx 1$, no correction --- while in the outer regimes it is stationary, $\sum_k\phi_k^{(1,3)}<1$, so deviations decay. The estimated band width is
\begin{equation}
\text{bw} = C_2 - C_1 .
\end{equation}
Dwyer et al.\ motivate \eqref{eq:tar} from the cost-of-carry model with non-zero transaction costs on minute-by-minute S\&P~500 futures/cash (1982--1990), and generalise it to a bivariate \emph{threshold error-correction mechanism}: with $y_t=(f_t,s_t)'$,
\begin{equation}
\Delta y_t = \nu^{(j)} + \alpha^{(j)} z_{t-1} + \sum_{i=1}^{p}\Gamma_i^{(j)}\Delta y_{t-i} + \varepsilon_t, \label{eq:tecm}
\end{equation}
where the error-correction loadings $\alpha^{(j)}=(\alpha_f^{(j)},\alpha_s^{(j)})'$ are regime-specific: $\alpha^{(2)}\approx 0$ in the band, $\alpha^{(1,3)}\ne 0$ outside. The two-equation form lets futures and cash adjust at \emph{different speeds}, which also identifies \emph{which market moves first} (price discovery).

\subsection*{Estimation}
Given $(C_1,C_2,d,p)$ the model is linear within regimes, so it is estimated by OLS regime-by-regime. The thresholds and delay are chosen by a grid search that minimises the residual sum of squares / an information criterion, subject to a minimum fraction of observations in each regime (Aggarwal imposes $\ge 40\%$ in the no-arbitrage regime). Lag order $p$ is set by AIC. Because the objective is discontinuous in $(C_1,C_2,d)$, the search is over a grid rather than by gradient methods.

\begin{takeaway}
Equation \eqref{eq:tar} is the exact model Aggarwal (2015) applied to Indian single-stock futures (\S\ref{sec:aggarwal}); it is our baseline band estimator. Its regime structure operationalises the theory: the middle regime \emph{is} the GV/DSSW no-arbitrage band, and $C_1,C_2$ are the (possibly asymmetric) friction wedges of \eqref{eq:gv-asym}.
\end{takeaway}

% =====================================================================
\section{Tse (2001): the smooth-transition (STAR) alternative}
\label{sec:tse}
A sharp threshold assumes every arbitrageur faces the \emph{same} cost. Tse (2001) argues that with heterogeneous arbitrageurs (different costs, capital, objectives) the \emph{aggregate} of many individual thresholds is smooth, not a step. He models the basis with a smooth-transition error-correction model in which the speed of adjustment rises \emph{continuously} with the size of the mispricing:
\begin{equation}
\Delta z_t = \big[\,\alpha_{\text{in}} + (\alpha_{\text{out}}-\alpha_{\text{in}})\,G(z_{t-d};\theta)\,\big]\,z_{t-1} + \sum_{i=1}^{p}\phi_i\,\Delta z_{t-i} + \varepsilon_t, \label{eq:star}
\end{equation}
where $G(\cdot)\in[0,1]$ is a transition function. Tse uses the \emph{exponential} form (ESTAR),
\begin{equation}
G(z_{t-d};\gamma,c) = 1 - \exp\!\big(-\gamma\,(z_{t-d}-c)^2\big), \qquad \gamma>0, \label{eq:estar}
\end{equation}
so that near equilibrium ($z_{t-d}\approx c$) $G\approx 0$ and there is little correction (the ``inside-band'' behaviour $\alpha_{\text{in}}\approx 0$), while for large deviations $G\to 1$ and adjustment occurs at the fast ``outside'' speed $\alpha_{\text{out}}<0$. The logistic form (LSTAR), $G=[1+\exp(-\gamma(z_{t-d}-c))]^{-1}$, allows \emph{asymmetric} adjustment for positive vs negative deviations --- directly relevant given the short-side asymmetry in India. Estimation is by nonlinear least squares / maximum likelihood, after a Teräsvirta LM test for the presence and shape (exponential vs logistic) of the non-linearity.

\begin{takeaway}
STAR is our robustness alternative to the ``abrupt'' TAR. If arbitrageurs are heterogeneous (they are --- proprietary desks, FPIs, and retail all face different costs), the smooth model fits better, and the LSTAR form gives a second, independent read on the band's asymmetry. Reporting both TAR and STAR guards against a knife-edge dependence on the threshold assumption.
\end{takeaway}

% =====================================================================
\section{Aggarwal (2015): the estimator applied to India, plus PCA commonality}
\label{sec:aggarwal-method}
Aggarwal's estimator is \eqref{eq:tar} run \emph{per stock, per day} at one-second frequency on 141 NSE single-stock futures/spot pairs, with the proportional basis \eqref{eq:basis-prop}, delay $d$ selected from a grid $\{1,5,\dots,300\}$ seconds and lag by AIC. The novel second step is to ask whether the estimated daily bandwidths $\{\text{bw}_{i,t}\}$ share a \emph{common} component. Stack them into a stock-by-day matrix, standardise each stock's series to zero mean and unit variance, and take the eigendecomposition of the sample covariance $\Sigma = \Gamma\Lambda\Gamma'$. The first principal component
\begin{equation}
\mathrm{PC1}_t = \gamma_1' \,\mathbf{bw}_t, \qquad \text{share of variance} = \frac{\lambda_1}{\sum_{k}\lambda_k}=\frac{\lambda_1}{N}, \label{eq:pca}
\end{equation}
is interpreted as the \emph{systematic} arbitrage-constraint factor. Its determinants are then examined by time-series regression on funding-liquidity and market variables (\S\ref{sec:aggarwal} reports the results).

\begin{takeaway}
This two-step design --- estimate per-stock bands, then PCA them into a systematic factor --- is precisely the skeleton of our Chapter~1 \emph{arbitrage-constraint index}. We extend it by (i) adding MWPL-specific inputs (ban dummy, \%MWPL utilisation, delivery share, ownership breadth) to the cross-section before extracting the common factor, and (ii) replacing Aggarwal's funding-liquidity driver with the \emph{ban / MWPL} shock as the exogenous mover of the band.
\end{takeaway}

% =====================================================================
\section{Hansen--Seo (2002) and Ters--Urban (2020): the TVECM we will use}
\label{sec:tvecm}
The estimators above either fix the cointegrating vector at $(1,-1)$ or treat the basis as univariate. Hansen \& Seo (2002) provide the modern, fully bivariate treatment: a \emph{two-regime threshold VECM with an unknown cointegrating vector} $\beta$. With $y_t=(f_t,s_t)'$ and error-correction term $w_{t-1}(\beta)=\beta' y_{t-1}$,
\begin{equation}
\Delta y_t = \begin{cases} A_1'\,X_{t-1}(\beta) + \varepsilon_t, & w_{t-1}(\beta)\le \hat\gamma,\\[3pt] A_2'\,X_{t-1}(\beta) + \varepsilon_t, & w_{t-1}(\beta) > \hat\gamma, \end{cases} \qquad X_{t-1}(\beta)=\big(1,\ w_{t-1}(\beta),\ \Delta y_{t-1}',\dots,\Delta y_{t-p}'\big)', \label{eq:hansenseo}
\end{equation}
estimated by maximum likelihood, jointly searching over the threshold $\gamma$ and $\beta$ on a two-dimensional grid; they also give a $\sup$-LM test of linear cointegration against the threshold alternative.

\paragraph{Ters--Urban (2020).} Ters and Urban extend Hansen--Seo to \emph{three} regimes (two thresholds $\gamma_1<\gamma_2$), add a constant in the cointegration relation (to absorb a persistent non-zero basis), and solve the harder joint estimation with a dynamic grid:
\begin{equation}
\Delta y_t = \begin{cases} A_1' X_{t-1}(\beta) + \varepsilon_t, & w_{t-1}(\beta) \le \gamma_1 \quad(\text{negative-basis arbitrage}),\\[3pt] A_2' X_{t-1}(\beta) + \varepsilon_t, & \gamma_1 < w_{t-1}(\beta) < \gamma_2 \quad(\text{no-arbitrage band}),\\[3pt] A_3' X_{t-1}(\beta) + \varepsilon_t, & w_{t-1}(\beta) \ge \gamma_2 \quad(\text{positive-basis arbitrage}). \end{cases} \label{eq:tersurban}
\end{equation}
The middle regime has (near-)zero error-correction loading; the outer regimes correct. The estimated thresholds recover the band, and its (possibly asymmetric) half-widths measure the \emph{overall round-trip arbitrage cost} on each side --- even when that cost is otherwise unobservable. Two properties make it ideal for us. First, the thresholds $\gamma_1,\gamma_2$ need not be symmetric, so the model \emph{estimates} the short-side asymmetry rather than assuming it. Second, and decisively, Ters and Urban design the method explicitly for \emph{pre/post-event analysis of regulatory changes}: comparing the estimated band before and after a dated shock measures how the policy moved arbitrage costs. That is exactly our 2024--26 reform design.

\begin{takeaway}
Ters--Urban \eqref{eq:tersurban} is the project's \emph{workhorse} band estimator (Chapter~4). We estimate the band on each stock's basis, allow the thresholds to depend on the Chapter~1 constraint index, and re-estimate around each dated reform (the 1~Oct~2025 MWPL/ban recast, the FutEq redefinition, the intraday-monitoring date) to read off how the band moved. The three-regime, asymmetric-threshold structure is the econometric embodiment of the one-sided GV wedge, and the ``unknown $\beta$ / constant in cointegration'' features handle India's persistent non-zero basis. TAR \eqref{eq:tar} and STAR \eqref{eq:star} are reported alongside as robustness.
\end{takeaway}

\part{Evidence from the Indian Market and the Institutional Mechanism}
\label{part:india}

This part grounds the theory and econometrics in Indian data and, crucially, documents the institutional mechanism the project exploits. It answers: what do Indian mispricings actually look like (magnitude, sign, asymmetry)? and what exactly is the MWPL ``F\&O ban,'' how is it triggered, and why is it a clean, exogenous shock to arbitrage capacity?

% =====================================================================
\section{Aggarwal (2015): the closest precedent and its results}
\label{sec:aggarwal}
Using the estimator of \S\ref{sec:aggarwal-method} on 141 NSE single-stock futures/spot pairs at one-second frequency (Mar--Aug 2013), Aggarwal reports:
\begin{itemize}[leftmargin=1.4em,itemsep=2pt]
\item \textbf{Asymmetric bands.} The average upper threshold is $\approx +3$ basis points; the average lower threshold is $\approx -15$ basis points. The no-arbitrage band is thus $\approx 18$ bps wide and \emph{strongly asymmetric}: the market tolerates much larger negative mispricing (spot rich / futures cheap) before correction, because reverse cash-and-carry requires shorting the stock, which is constrained. This is \eqref{eq:gv-asym} made visible in Indian data.
\item \textbf{Cross-sectional variation.} Band width and asymmetry both increase as stock liquidity falls (from the most- to least-liquid quartile) --- consistent with liquidity/transaction costs setting the wedge.
\item \textbf{Commonality.} By the PCA of \eqref{eq:pca}, the first principal component explains $\approx 17\%$ of the cross-sectional variation in bandwidths (first three PCs $\approx 37\%$): there is a systematic arbitrage-constraint factor.
\item \textbf{Funding liquidity drives the common factor.} Regressing $\mathrm{PC1}_t$ on candidate determinants, the single strongest is funding liquidity (proxied by MIBOR), with $R^2\approx 18\%$. In a natural experiment --- the RBI's mid-July 2013 liquidity tightening (overnight MIBOR jumped from $7.88\%$ to $9.23\%$) --- the systematic band component spiked, an ``unintended consequence of an action by the regulator.''
\end{itemize}

\begin{takeaway}
Aggarwal is the paper we most directly build on, and the gap is precise. She \emph{measures} Indian no-arbitrage bands and links their common component to \emph{funding liquidity}; she does \emph{not} use the F\&O ban / MWPL mechanism, and she establishes no causal effect of a \emph{constraint shock} on the band. Our contribution $=$ her band-measurement method $+$ the MWPL ban and the 2024--26 rule changes as the exogenous mover of the band $+$ a decade-long ($2015$--$2026$), daily, all-F\&O-stocks panel in place of her six-month, 141-stock, intraday sample. Her $-15/+3$ bp asymmetry is the number our estimates must first reproduce, then explain.
\end{takeaway}

% =====================================================================
\section{Nandan et al.\ (2014), Roy \& Chakraborty (2023), Slivka et al.\ (2011)}
\label{sec:india-basis}

\paragraph{Nandan, Agrawal \& Bhargava (2014) --- the basis is persistently negative.} Using the cost-of-carry model \eqref{eq:coc} on Nifty index futures (daily, 2003--2012; $2{,}382$ observations after filtering expiry days; $r=$ 30-day MIBOR), they find mean mispricing of $-\rupee 8.19$, with $71\%$ of observations \emph{negative} (futures under-priced) versus $29\%$ positive, and the negative side larger in magnitude ($-\rupee 13.99$ vs $+\rupee 6.04$). Non-parametric tests (Kolmogorov--Smirnov for non-normality; Mann--Whitney across the sign of the basis; Kruskal--Wallis across days-to-expiry) confirm the mispricing is significant, rises with days-to-expiry, and \emph{attenuated over the decade} (mean $\to -\rupee 0.23$ by 2012). The persistent negativity is the index-level signature of the short-side constraint --- and the time trend tells us the pre-2015 market behaves differently, justifying a regime-aware sample.

\paragraph{Roy \& Chakraborty (2023) --- cointegration and price discovery.} Using Johansen cointegration and lead--lag/causality tests on NSE single-stock futures and their underlying, they confirm that spot and futures are cointegrated (the law of one price holds in the \emph{long run}) with short-run deviations and a price-discovery role for the futures segment. This is the empirical warrant for the cointegration assumption underlying the TVECM of \S\ref{sec:tvecm}: the basis is a stationary error-correction term to which prices eventually return.

\paragraph{Slivka, Shah \& Wu (2011) --- the band is about twice as wide as in mature markets.} This practitioner study finds that the zero-arbitrage band for single-stock-vs-futures arbitrage in India is ``typically at least twice that occurring in mature global markets,'' attributing the width to high brokerage/exchange/regulatory fees and the securities-transaction tax, an under-developed securities-lending market, and execution frictions. It reports that real arbitrage is confined almost entirely to \emph{intraday single-stock-vs-futures} trades, with index-basket arbitrage largely absent.

\begin{takeaway}
Three usable facts: (i) the Indian basis is persistently negative and asymmetric (Nandan) --- our sign predictions are grounded; (ii) spot and futures are cointegrated (Roy--Chakraborty) --- the TVECM is valid; (iii) Indian bands run $\approx 2\times$ mature-market width (Slivka) --- a concrete magnitude our estimates should match, and confirmation that single-stock futures (not index baskets, not cash short-sales) are the arbitrage channel to study.
\end{takeaway}

% =====================================================================
\section{The institutional mechanism: MWPL, the F\&O ban, and the 2024--26 reforms}
\label{sec:institution}

\subsection*{A correction on Indian short-selling (and why it strengthens the design)}
It is often said loosely that ``there is no short selling in India.'' The precise position matters for our identification and is this: \emph{intraday} short selling in the cash market is freely available; \emph{naked overnight} short selling is prohibited; and the Securities Lending and Borrowing (SLB) market that would support covered overnight shorts is thin and lightly used. The consequence is that the practical way to take a persistent short/negative-basis arbitrage position is through \emph{single-stock futures} --- which is exactly the instrument the MWPL ban restricts. Far from weakening the project, this sharpens it: the binding constraint sits on the very channel arbitrageurs must use, so a ban on fresh futures positions is a direct contraction of arbitrage capacity, not a peripheral one.

\subsection*{The ban mechanism}
Each F\&O stock has a \emph{Market-Wide Position Limit} (MWPL). When market-wide open interest in that stock crosses $95\%$ of its MWPL, the stock enters the \emph{ban period}: only position-\emph{reducing} trades are allowed, no fresh positions. The stock exits the ban when open interest falls back below $80\%$. NSE publishes the securities-in-ban list and each stock's MWPL and \%-utilisation \emph{daily}. The definition of MWPL changed with the 2025 reform:
\begin{align}
\text{Pre-reform (to 30 Sep 2025):}\quad & \mathrm{MWPL} = 20\%\ \text{of free-float shares} \ \ (\text{or } 30\times \text{prior-month avg.\ volume}), \label{eq:mwpl-old}\\
\text{From 1 Oct 2025:}\quad & \mathrm{MWPL} = \min\!\big(15\%\ \text{free float},\ 65\times \mathrm{ADDV}\big),\ \text{floor } 10\%\ \text{free float}, \label{eq:mwpl-new}
\end{align}
where ADDV is the \emph{Average Daily Delivery Value} in the cash market, and the new MWPL is \emph{recomputed quarterly} on a rolling three-month ADDV. Simultaneously, open interest itself was redefined from notional to \emph{future-equivalent (FutEq)}: each position is counted in delta-adjusted terms,
\begin{equation}
\text{FutEq OI} = \Big|\textstyle\sum_{\text{positions}} \Delta_{\text{position}}\times \text{notional}\Big|, \qquad \Delta_{\text{long fut}}=+1,\ \Delta_{\text{option}}\in[-1,+1], \label{eq:futeq}
\end{equation}
and the ban-period rule was recast so that any trade after entry must reduce \emph{end-of-day FutEq delta} (a passive increase from price moves is not a breach; a sign flip does not count as reduction).

\subsection*{The reform timeline (dated, market-wide shocks)}
The 2024--26 overhaul provides a dense sequence of exogenous, precisely dated changes to the constraints the project studies. Table~\ref{tab:reforms} collects them from the two SEBI circulars.

\begin{table}[h]
\centering\small
\renewcommand{\arraystretch}{1.25}
\begin{tabular}{@{}p{0.52\textwidth} p{0.20\textwidth} c@{}}
\toprule
\textbf{Reform} & \textbf{Effective} & \textbf{Source} \\
\midrule
Index contract size raised to \rupee15--20 lakh (from \rupee5--10 lakh) & 20 Nov 2024 & SEBI 2024/132 \\
Weekly expiries limited to one benchmark index per exchange & 20 Nov 2024 & SEBI 2024/132 \\
Additional 2\% ELM on short options on expiry day & 20 Nov 2024 & SEBI 2024/132 \\
Upfront collection of option premium from buyers & 1 Feb 2025 & SEBI 2024/132 \\
Calendar-spread margin benefit removed on expiry day & 1 Feb 2025 & SEBI 2024/132 \\
Intraday monitoring of index position limits ($\ge 4$ snapshots) & 1 Apr 2025 & SEBI 2024/132 \\
Index-futures/options position limits (net \rupee1{,}500 cr / gross \rupee10{,}000 cr) & 1 Jul 2025 & SEBI 2025/79 \\
FutEq (delta-based) OI in force & mid-2025 & SEBI 2025/79 \\
\textbf{New MWPL formula \eqref{eq:mwpl-new} + delta-based ban-period rule} & \textbf{1 Oct 2025} & SEBI 2025/79 \\
Individual entity single-stock position limits recalibrated to new MWPL & 1 Oct 2025 & SEBI 2025/79 \\
Intraday MWPL monitoring for single stocks; $\ge 14$ constituents for non-benchmark indices & 3 Nov 2025 & SEBI 2025/79 \\
Pre-open session extended to F\&O; index-option limits to normal enforcement & 6 Dec 2025 & SEBI 2025/79 \\
\bottomrule
\end{tabular}
\caption{The 2024--26 SEBI derivatives reforms, with exact effective dates. One stated objective of the May 2025 circular is explicitly ``to reduce instances of spurious F\&O ban periods in single stocks'' --- i.e.\ SEBI is directly re-engineering the mechanism this project studies.}
\label{tab:reforms}
\end{table}

\begin{takeaway}
This section defines the project's identifying variation. The ban is \emph{mechanical} (a bright-line 95\% rule), \emph{dated} (published daily), and \emph{repeated} (hundreds of stock-episodes). Two features are decisive for causal inference (developed in Part~VI): (1) the MWPL \emph{denominator} in \eqref{eq:mwpl-old}--\eqref{eq:mwpl-new} moves for reasons unrelated to today's speculation --- free float changes with issuance, and post-reform the binding term is a \emph{quarterly-reset, trailing} delivery-value figure (ADDV) --- so the same open interest can be in- or out-of-ban depending on stale, exogenous inputs, giving an instrument for ban status; and (2) the dated reforms in Table~\ref{tab:reforms} are market-wide shocks to constraint \emph{intensity} that the Ters--Urban TVECM (\S\ref{sec:tvecm}) can read off in a pre/post design. The FutEq redefinition \eqref{eq:futeq} also tells us the post-2025 constraint variable must be built in delta terms, which requires option deltas from settlement-implied volatilities.
\end{takeaway}

\part{The Behavioural / Retail Demand Side}
\label{part:behavioural}

The arbitrage core (Parts~I--IV) is the \emph{supply} side of price correction --- who fixes mispricings, and what stops them. This part is the \emph{demand} side --- who \emph{creates} the mispricing pressure, and why it concentrates in exactly the stocks with the widest bands. In DSSW's language (\S\ref{sec:dssw}) this is the noise-trader block; in our project it is the concluding policy chapter, which argues that retail crowding is a leading reason single stocks hit the 95\% MWPL trigger in the first place. It requires no new firm-level retail data (which public sources cannot provide) --- it borrows the identified results below.

% =====================================================================
\section{Kumar (2009): who gambles, and in which stocks}
\label{sec:kumar}
Kumar shows individual investors' propensity to gamble spills into their portfolios: they over-weight \emph{lottery-type stocks}. A lottery stock is cheap, volatile, and positively skewed --- a small outlay with a small chance of a large payoff. Operationally, Kumar flags a stock as lottery-type when it is simultaneously
\begin{equation}
\text{low price}, \qquad \text{high idiosyncratic volatility } \sigma(\hat\varepsilon_i), \qquad \text{high idiosyncratic skewness } \mathrm{skew}(\hat\varepsilon_i), \label{eq:lottery}
\end{equation}
where $\hat\varepsilon_i$ are residuals from a factor model (e.g.\ Fama--French). He documents that lottery demand rises among younger, poorer, less-educated investors and in economic downturns, and that lottery stocks subsequently \emph{under-perform} --- over-demand pushes prices up and expected returns down.

\begin{takeaway}
Equation~\eqref{eq:lottery} identifies \emph{which} stocks retail crowds into: low-priced, high-IVOL, high-skew names. These are the same characteristics that, in Part~I, produce wide no-arbitrage bands (high idiosyncratic risk $\Rightarrow$ large $\sigma_A^2(1-\rho^2)$ in \eqref{eq:gv-price}). Demand and constraint therefore line up on the same stocks --- the mechanism behind hypothesis H3 (bands widest for small, high-IVOL, retail-heavy names) and the demand-side reason such names approach the MWPL trigger.
\end{takeaway}

% =====================================================================
\section{Sehgal, Vasishth \& Deisting (2024): the lottery factor in India}
\label{sec:lottery-india}
This paper confirms Kumar's mechanism for India (Dec 2001--Mar 2021) and builds a tradable factor. Following Bali, Cakici \& Whitelaw (2011), the key lottery proxy is \emph{MAX}: the average of the $k$ largest daily returns of a stock in a month,
\begin{equation}
\mathrm{MAX}^{(k)}_{i,t} = \frac{1}{k}\sum_{\ell=1}^{k}\,R^{(\ell)}_{i,t}, \label{eq:max}
\end{equation}
where $R^{(1)}_{i,t}\ge R^{(2)}_{i,t}\ge\cdots$ are the ranked daily returns within month $t$ (typically $k=1$ to $5$). Sorting on MAX (alongside skewness, tail risk, and IVOL) and forming a long-low/short-high portfolio yields a \emph{lottery factor}; augmenting the Fama--French five-factor model with it,
\begin{equation}
r_{i,t}-r_{f,t} = \alpha_i + \beta_i^{\text{MKT}}\mathrm{MKT}_t + \beta_i^{\text{SMB}}\mathrm{SMB}_t + \beta_i^{\text{HML}}\mathrm{HML}_t + \beta_i^{\text{RMW}}\mathrm{RMW}_t + \beta_i^{\text{CMA}}\mathrm{CMA}_t + \beta_i^{\text{LOT}}\mathrm{LOT}_t + \varepsilon_{i,t}, \label{eq:ff6}
\end{equation}
produces a better-performing (six-factor) benchmark. They find MAX predicts returns beyond other risk measures, and that lottery behaviour in India is \emph{mainly retail-driven} and tied to probability mis-weighting and overreaction to attention-grabbing events.

\begin{takeaway}
The Indian lottery effect already exists in the literature, so we \emph{cite rather than re-derive} it. It supplies the demand-side counterpart to our supply-side story and shares its toolkit (MAX, IVOL, skew) with our constraint index. It also warns against re-treading crowded ground --- our contribution is the constraint/ban mechanism, not the lottery factor.
\end{takeaway}

% =====================================================================
\section{Barber \& Odean (2000, 2013): retail investors as net noise}
\label{sec:barber-odean}
The two Barber--Odean references establish that retail investors, in aggregate, trade too much and lose. In ``Trading Is Hazardous to Your Wealth'' (2000), using $66{,}465$ US discount-brokerage households (1991--96), they show net returns fall with turnover: the highest-turnover quintile earns an annual net return of $\approx 11.4\%$ against a market return of $\approx 17.9\%$, a gap driven by trading costs and overconfident over-trading. The survey ``The Behavior of Individual Investors'' (2013) synthesises the field: individual investors under-perform, hold under-diversified portfolios, exhibit the disposition effect (selling winners, holding losers), and buy attention-grabbing / past-return stocks --- i.e.\ they behave as DSSW noise traders.

\begin{takeaway}
This is the micro-foundation for treating heavy retail activity as noise-trader demand rather than informed order flow, licensing the DSSW interpretation of retail crowding in the policy chapter. It is US evidence; we use it for framing, not identification.
\end{takeaway}

% =====================================================================
\section{SEBI (2024): the aggregate retail-loss study}
\label{sec:sebi-loss}
SEBI's study of individual traders in equity F\&O (FY22--FY24) is the authoritative, account-level evidence that public data cannot reproduce:
\begin{itemize}[leftmargin=1.4em,itemsep=1pt]
\item $93\%$ of over 1 crore individual F\&O traders lost money over FY22--24; average loss $\approx \rupee 2$ lakh per trader; only $1\%$ earned more than $\rupee 1$ lakh.
\item In FY24, proprietary traders and FPIs booked gross profits of $\approx \rupee 33{,}000$ cr and $\approx \rupee 28{,}000$ cr respectively, while individuals lost $\approx \rupee 61{,}000$ cr; $96\%$--$97\%$ of the prop/FPI profits came from \emph{algorithmic} trading.
\item The base skews young and less wealthy: under-30 traders rose from $31\%$ (FY23) to $43\%$ (FY24); $\approx 72\%$ from beyond the top-30 cities; $>75\%$ declared income below $\rupee 5$ lakh; and $>75\%$ of loss-makers kept trading.
\end{itemize}

\begin{takeaway}
This is the empirical face of DSSW's relative-returns result (\S\ref{sec:dssw}): noise traders (retail) transfer wealth to sophisticated players (algorithmic prop/FPI desks) precisely by bearing the risk they create. It is the authoritative welfare number, which is exactly why the project does \emph{not} attempt to re-estimate retail welfare from public aggregates (we would produce a worse version of a number the regulator already owns). It motivates \emph{why} SEBI is tightening the very constraints we study --- closing the loop between the demand side and the policy relevance of Project~1.
\end{takeaway}

\part{Synthesis: What the Literature Implies for Our Model}
\label{part:synthesis}

The review has a single destination: a formal model of our project, to be built \emph{before} the empirics so that the hypotheses are theory-derived rather than data-fitted. This part draws the strands together, states what such a model must reproduce, and sketches the natural candidate.

% =====================================================================
\section{The through-line}
Every paper is a facet of one object: the \emph{asymmetric no-arbitrage band} around an Indian stock's cost-of-carry fair value, and how an exogenous contraction of arbitrage capacity widens it. Formally, write the observed (log) basis of stock $i$ as
\begin{equation}
z_{i,t} = \underbrace{\theta_{i,t}}_{\substack{\text{risk-bearing wedge}\\ \sim \alpha\sigma_i^2(1-\rho_i^2)u_{i,t}\ \text{(GV, DSSW)}}} \;\pm\; \underbrace{c_{i,t}}_{\substack{\text{friction wedge, one-sided}\\ \text{short-sale/position limit (GV)}}} \;+\; \underbrace{a_{i,t}}_{\substack{\text{amplification, state-dependent}\\ \text{PBA fire-sale (SV)}}} \;+\; \eta_{i,t},
\end{equation}
and let the \emph{band} be the interval $[\,C_1(\,\cdot\,),\,C_2(\,\cdot\,)\,]$ of $z$ within which no correction occurs (Part~III). The project's claims are then: (i) the band is asymmetric because $c$ is one-sided on the short/negative-basis leg (GV \eqref{eq:gv-asym}; Aggarwal $-15/+3$ bp; Ofek $13.6\%/4.4\%$); (ii) the band widens with a constraint index built from ban status, \%MWPL, idiosyncratic volatility, delivery share and ownership breadth (DSSW variance channel; Aggarwal PCA); and (iii) the MWPL ban is an exogenous jump in the position-limit component of $c$, so it should widen the band and slow correction --- causally, not just in correlation.

% =====================================================================
\section{Paper-to-project map}
Table~\ref{tab:map} states, for each reference, what the project takes and where it is used.

\begin{table}[h]
\centering\small
\renewcommand{\arraystretch}{1.2}
\begin{tabular}{@{}p{0.30\textwidth} p{0.42\textwidth} p{0.20\textwidth}@{}}
\toprule
\textbf{Paper} & \textbf{What we take} & \textbf{Used in} \\
\midrule
Shleifer--Vishny (1997) & Capital is state-dependent; shocks to capacity widen mispricing; fire-sale amplification & Motivation; Ch.\ 4 (non-linearity) \\
De Long et al.\ (1990) & Mispricing scales with variance of non-fundamental demand; systematic sentiment & Constraint index (IVOL); policy ch. \\
Gromb--Vayanos (2010) & Price $=$ risk term $+$ one-sided friction wedge; position limits as shadow cost & Core model; band asymmetry \\
Cost-of-carry basis & Primary mispricing measure $z=[F-Se^{(r-\delta)\tau}]/S$ & Ch.\ 2 (measure 1) \\
Ofek et al.\ (2004) & Put--call-parity deviation; short-side asymmetry; predictability & Ch.\ 2 (measure 2) \\
PEAD-India (2018) & SUE-based drift; Fama--MacBeth; robust to IVOL & Ch.\ 2 (measure 3) \\
Stambaugh et al.\ (2015) & Composite mispricing score; IVOL as arbitrage risk; Fama--MacBeth & Ch.\ 5 (cross-section) \\
Dwyer (1996); YPP (1994) & Three-regime TAR / threshold-ECM for the basis & Ch.\ 4 (baseline) \\
Tse (2001) & Smooth-transition (STAR) alternative; heterogeneity; asymmetry & Ch.\ 4 (robustness) \\
Aggarwal (2015) & Band estimation $+$ PCA commonality on Indian SSF; the precedent/gap & Ch.\ 1 \& 4 (method) \\
Hansen--Seo (2002); Ters--Urban (2020) & 3-regime TVECM, unknown $\beta$, pre/post-event design & Ch.\ 4 (workhorse) \\
Nandan (2014); Roy--Chakraborty (2023); Slivka (2011) & Negative/asymmetric basis; cointegration; $2\times$ band benchmark & Calibration \\
Kumar (2009); Sehgal et al.\ (2024) & Which stocks retail crowds (low-price/high-IVOL/high-skew) & Policy ch.; H3 \\
Barber--Odean (2000, 2013) & Retail as net noise traders & Policy ch. \\
SEBI (2024); SEBI circulars; NSE rules & Welfare number; the ban mechanism and dated shocks & Policy ch.; identification \\
\bottomrule
\end{tabular}
\caption{How each reference feeds the project.}
\label{tab:map}
\end{table}

% =====================================================================
\section{What the theoretical model must reproduce}
\label{sec:model-requirements}
The literature pins down a checklist for the next-stage model. A satisfactory model of our project should deliver, from primitives, all of the following --- which are then the \emph{pre-registered} hypotheses, derived not assumed:
\begin{enumerate}[leftmargin=1.5em,itemsep=3pt]
\item \textbf{A no-arbitrage band of explicit width.} An interval $[C_1,C_2]$ in basis space with a closed-form (or comparative-static) expression for its half-widths in terms of transaction costs, risk-bearing capacity, and a \emph{position-limit shadow cost}. (Maps to H1.)
\item \textbf{Asymmetry.} $\abs{C_1} > \abs{C_2}$ (wider on the short/negative-basis side) whenever the binding constraint is on shorting/short-side futures --- the GV \eqref{eq:gv-asym} property, matching Aggarwal's $-15/+3$ bp.
\item \textbf{Monotonic comparative static in the constraint.} $\partial(C_2-C_1)/\partial(\text{constraint}) > 0$: tighter constraints widen the band. The MWPL ban is the discrete limiting case (shadow price jumps when the limit binds). (Maps to H1, H4.)
\item \textbf{Slower correction inside a wider band.} A longer mispricing half-life when the band is wider / the stock is in the ban --- the dynamic, testable via the TVECM error-correction speed and via PEAD depth. (Maps to H2.)
\item \textbf{A cross-sectional pricing implication.} Signals/anomalies earn higher returns in high-constraint stocks (mispricing not risk) --- SYY logic under India's near-absolute short constraint. (Maps to H5.)
\item \textbf{A falsification route.} The band should respond to \emph{exogenous} movements in the constraint (the ADDV/free-float denominator revisions; the dated reforms), not merely to endogenous crowding. (Maps to H3/H6.)
\end{enumerate}

% =====================================================================
\section{A candidate model skeleton for the next step}
\label{sec:skeleton}
The most direct route is to take the Gromb--Vayanos arbitrageur (\S\ref{sec:gv}) and add a \emph{hard position limit} of MWPL type. Sketch: a risk-averse arbitrageur chooses a futures position $x$ against a demand shock $u$, maximising the mean--variance objective \eqref{eq:gv-mv} subject to a position cap
\begin{equation}
\abs{x} \le L, \qquad L = \kappa\cdot\mathrm{MWPL} = \kappa\cdot\min\!\big(0.15\,\mathrm{FF},\ 65\,\mathrm{ADDV}\big),
\end{equation}
representing the per-participant share of the market-wide limit. The Karush--Kuhn--Tucker conditions introduce a multiplier $\lambda\ge 0$ on the cap; the equilibrium basis becomes
\begin{equation}
z = \alpha\sigma^2(1-\rho^2)\,u \;\pm\; c \;+\; \lambda(u,L)\cdot\mathrm{sgn}(x), \qquad \lambda>0 \iff \abs{x}=L .
\end{equation}
When market-wide OI approaches $95\%$ of MWPL, the cap binds ($\lambda$ jumps from $0$), the wedge widens discretely, and the band's constrained-side half-width increases --- delivering requirements (1)--(4) analytically. Because $L$ depends on ADDV/free-float through \eqref{eq:mwpl-new}, exogenous denominator revisions shift $L$ (and hence $\lambda$) independently of current mispricing, delivering requirement (6). Embedding this in a two-period (or infinite-horizon, à la DSSW) setting with stochastic $u$ generates the persistence dynamics of requirement (3); adding the SV fund-flow channel \eqref{eq:sv-G} yields the amplification of \eqref{eq:sv-resiliency}. This is a sketch, not the model --- but it shows the literature furnishes every block needed, which is the point of building the theory first.

% =====================================================================
\section{Where the gap sits (our contribution)}
The international theory of limits to arbitrage (SV, DSSW, GV) is mature; the measurement (basis, PCP, PEAD, SYY) and the band econometrics (TAR/STAR/TVECM) are standard; and the Indian evidence (Aggarwal, Nandan, Roy--Chakraborty, Slivka) establishes that wide, asymmetric, cointegrated bands exist. \emph{No paper ties a mechanical, dated position-limit shock (the MWPL ban) to the width and dynamics of the Indian no-arbitrage band, nor uses it as an instrument for arbitrage capacity.} That junction --- theory-modelled first, then tested on a decade of public NSE data around the 2024--26 reforms --- is the project's contribution. The immediate next step is to turn \S\ref{sec:model-requirements}--\S\ref{sec:skeleton} into a fully worked model with proofs, then take its derived predictions to the data.


\newpage
\phantomsection
\addcontentsline{toc}{part}{References}
\begin{thebibliography}{99}
\setlength{\itemsep}{2pt}

\bibitem{aggarwal2015} Aggarwal, N.\ (2015). \emph{Limits to Arbitrage: The Case of Single Stock Futures and Spot Prices.} IGIDR Working Paper WP-2015-010, Indira Gandhi Institute of Development Research, Mumbai.

\bibitem{balkefomby1997} Balke, N.\ S., \& Fomby, T.\ B.\ (1997). Threshold Cointegration. \emph{International Economic Review}, 38(3), 627--645.

\bibitem{bali2011} Bali, T.\ G., Cakici, N., \& Whitelaw, R.\ F.\ (2011). Maxing Out: Stocks as Lotteries and the Cross-Section of Expected Returns. \emph{Journal of Financial Economics}, 99(2), 427--446.

\bibitem{barberodean2000} Barber, B.\ M., \& Odean, T.\ (2000). Trading Is Hazardous to Your Wealth: The Common Stock Investment Performance of Individual Investors. \emph{Journal of Finance}, 55(2), 773--806.

\bibitem{barberodean2013} Barber, B.\ M., \& Odean, T.\ (2013). The Behavior of Individual Investors. In \emph{Handbook of the Economics of Finance}, Vol.\ 2B, 1533--1570. Elsevier.

\bibitem{delong1990} De Long, J.\ B., Shleifer, A., Summers, L.\ H., \& Waldmann, R.\ J.\ (1990). Noise Trader Risk in Financial Markets. \emph{Journal of Political Economy}, 98(4), 703--738.

\bibitem{dwyer1996} Dwyer, G.\ P., Locke, P., \& Yu, W.\ (1996). Index Arbitrage and Nonlinear Dynamics Between the S\&P 500 Futures and Cash. \emph{Review of Financial Studies}, 9(1), 301--332.

\bibitem{gromb2010} Gromb, D., \& Vayanos, D.\ (2010). Limits of Arbitrage: The State of the Theory. \emph{Annual Review of Financial Economics}, 2, 251--275. (NBER Working Paper 15821.)

\bibitem{hansenseo2002} Hansen, B.\ E., \& Seo, B.\ (2002). Testing for Two-Regime Threshold Cointegration in Vector Error-Correction Models. \emph{Journal of Econometrics}, 110(2), 293--318.

\bibitem{harshita2018} Harshita, Singh, S., \& Yadav, S.\ S.\ (2018). Post-Earnings-Announcement Drift Anomaly in India: A Test of Market Efficiency. \emph{Theoretical Economics Letters}, 8(14), 3178--3195.

\bibitem{kumar2009} Kumar, A.\ (2009). Who Gambles in the Stock Market? \emph{Journal of Finance}, 64(4), 1889--1933.

\bibitem{nandan2014} Nandan, T., Agrawal, P., \& Bhargava, S.\ (2014). Mispricing in CNX Nifty Futures: An Empirical Investigation. \emph{Asia-Pacific Journal of Management Research and Innovation}, 10(4), 413--422.

\bibitem{ofek2004} Ofek, E., Richardson, M., \& Whitelaw, R.\ F.\ (2004). Limited Arbitrage and Short Sales Restrictions: Evidence from the Options Markets. \emph{Journal of Financial Economics}, 74(2), 305--342. (NBER Working Paper 9423.)

\bibitem{roy2023} Roy, P.\ S., \& Chakraborty, T.\ (2023). Efficiency of Indian Equity Futures Market --- An Empirical Analysis with Reference to the National Stock Exchange. \emph{Global Business Review}, 24(6), 1326--1352.

\bibitem{sebi2024study} Securities and Exchange Board of India (2024). \emph{Study on Profits and Losses of Individual Traders in the Equity F\&O Segment} (PR No.\ 22/2024, September 2024).

\bibitem{sebi2024circular} Securities and Exchange Board of India (2024). \emph{Measures to Strengthen Equity Index Derivatives Framework for Increased Investor Protection and Market Stability.} Circular SEBI/HO/MRD/TPD-1/P/CIR/2024/132, 1 October 2024.

\bibitem{sebi2025circular} Securities and Exchange Board of India (2025). \emph{Measures for Enhancing Trading Convenience and Strengthening Risk Monitoring in Equity Derivatives.} Circular SEBI/HO/MRD/TPD-1/P/CIR/2025/79, 29 May 2025.

\bibitem{sehgal2024} Sehgal, S., Vasishth, V., \& Deisting, F.\ (2024). Lottery Factor and Stock Returns: Evidence from India. \emph{Borsa Istanbul Review}, 24(3), 449--459.

\bibitem{shleifervishny1997} Shleifer, A., \& Vishny, R.\ W.\ (1997). The Limits of Arbitrage. \emph{Journal of Finance}, 52(1), 35--55. (NBER Working Paper 5167.)

\bibitem{slivka2011} Slivka, R.\ T., Shah, V., \& Wu, J.\ (2011). \emph{Arbitrage of Single Stocks versus Futures in India.} Working Paper, Polytechnic Institute of New York University.

\bibitem{stambaugh2015} Stambaugh, R.\ F., Yu, J., \& Yuan, Y.\ (2015). Arbitrage Asymmetry and the Idiosyncratic Volatility Puzzle. \emph{Journal of Finance}, 70(5), 1903--1948. (NBER Working Paper 18560.)

\bibitem{tersurban2020} Ters, K., \& Urban, J.\ (2020). Estimating Unknown Arbitrage Costs: Evidence from a 3-Regime Threshold Vector Error Correction Model. \emph{Journal of Financial Markets}, 47, 100503.

\bibitem{tse2001} Tse, Y.\ (2001). Index Arbitrage with Heterogeneous Investors: A Smooth Transition Error Correction Analysis. \emph{Journal of Banking \& Finance}, 25(10), 1829--1855.

\bibitem{yadav1994} Yadav, P.\ K., Pope, P.\ F., \& Paudyal, K.\ (1994). Threshold Autoregressive Modeling in Finance: The Price Differences of Equivalent Assets. \emph{Mathematical Finance}, 4(2), 205--221.

\bibitem{nse_mwpl} National Stock Exchange of India. \emph{Market-Wide Position Limits and F\&O Ban Methodology; Daily Reports on Securities in Ban Period and MWPL Utilisation.} NSE, Mumbai.

\end{thebibliography}


\end{document}
